\chapter{1932:Irving Langmuir}

\label{chap:chemistry1932}

The Nobel Prize in Chemistry for the year 1932 was awarded to Irving Langmuir.

\textbf{Motivation:}

for his discoveries and investigations in surface chemistry

\section{Irving Langmuir}

\subsection*{Early Life and Education}

Irving Langmuir was born on 1881-01-31 in Brooklyn, New York, United States.
He was an American physical chemist awarded the Nobel Prize in Chemistry in 1932 for his discoveries and investigations in surface chemistry.

\subsection*{Career and Major Research}

Langmuir studied at Columbia University and received his doctorate at the University of Göttingen in 1906 under Walther Nernst.
From 1909 he worked at the General Electric Research Laboratory in Schenectady, New York, where he remained until his retirement in 1950.
His early work at General Electric concerned the incandescent lamp; he developed gas-filled lamps with coiled tungsten filaments, greatly improving their efficiency and lifetime.
He investigated the adsorption of gases on surfaces and the properties of monomolecular films on liquids, laying the foundations of modern surface chemistry.
He also worked on thermionic emission, on the vacuum pump, and on the atomic-hydrogen welding process, and he introduced the concept of electron shells in atoms.
Later in life he studied the artificial production of rain.

\subsection*{Nobel Prize-winning Work}

The work for which Irving Langmuir was awarded the Nobel Prize in Chemistry in 1932 was described by the Nobel Committee as: "for his discoveries and investigations in surface chemistry".

Langmuir's prize recognized his discoveries concerning the behavior of surfaces.
He established quantitative laws for the adsorption of gases on solids, showed that certain films spread on water are exactly one molecule thick, and developed a general picture of chemical reactions at surfaces.

\subsection*{Significance and Impact}

Langmuir's work created the scientific discipline of surface chemistry and provided the basis for understanding catalysis on solid surfaces, lubrication, and the behavior of thin films.
The Langmuir adsorption isotherm and the journal Langmuir are named in his honour.

\subsection*{Later Career and Honors}

Langmuir remained at the General Electric Research Laboratory until 1950.
He received the Franklin Medal, was elected a foreign member of the Royal Society, and held honorary degrees from several universities.
He died on 1957-08-16 in Falmouth, Massachusetts, United States.

\subsection*{Legacy}

The legacy of Irving Langmuir endures through the lasting impact of his scientific contributions.
Surface chemistry, the discipline he founded, remains central to catalysis, electronics, and materials science.

\subsection*{Selected Publications}

\begin{itemize}

\item Langmuir, I. (1916--1918). "The Constitution and Fundamental Properties of Solids and Liquids." Journal of the American Chemical Society.

\item Langmuir, I. (1919). "The Arrangement of Electrons in Atoms and Molecules." Journal of the American Chemical Society.

\item Langmuir, I. (1932). "Surface Chemistry." Nobel Lecture, Stockholm.

\end{itemize}

\clearpage

\section*{Chapter Summary}

The 1932 prize recognized Irving Langmuir for his discoveries and investigations in surface chemistry. His studies of adsorption, surface films, and the monomolecular layer laid the foundations of modern surface science.
